\documentclass{article}
\usepackage{pathfinder-readable}

\title{A Finite Prize for Coordinating Agents}

\begin{document}
\maketitle

\begin{abstract}
This account connects a paper on incentives for AI agents with a simulation in which agents spend
energy to work and survive. The thread asked whether changing the division of an existing job prize
could improve job choice without adding energy. The peers found an exact incentive problem in the
simulation's split prize and a bounded priority rule that makes a group-energy optimum a one-round
Nash equilibrium under fixed, known success and cost laws. The result does not show that the rule
improves the simulation's measured energy or survival across rounds. The verifier therefore left the
thread at DRAFT, with a controlled comparison still needed.
\end{abstract}

\section{Introduction}

\emph{Mechanism Design for Alignment and Control}, called Q here, studies how to make AI agents
report information and take desired actions when their preferences and abilities are uncertain. Its
peer scores and coupled rewards can make one agent's payoff depend on another's action. Some
constructions scale rewards freely or use an unbounded penalty. Their cost would matter if rewards
had to be paid as a scarce resource \cite{Q}.

\emph{The Energy Society}, called P here, studies agents in a repeated simulation. Generating tokens
costs energy, successful jobs replenish it, and an agent with no energy deactivates. Agents can
recommend jobs, but each chooses its own action. Agents who solve the same job split one fixed prize
equally. The paper reports job choices and collisions alongside activity and donations \cite{P}.

The connexion pursued was Q's idea of coupled incentives applied to P's job choice. A recommendation
alone does not oblige an agent to follow it, while P's shared prize may reward joining a job that
another agent already attempts. The peers asked whether a rule could change that incentive by
reallocating the existing prize. They treated energy as a real budget and separated the question of
one-round job choice from the harder questions of private information and survival over many rounds.

\section{What was done}

The peers first compared P's private prize with the group's payout. A second solver could be paid
even when the first had already earned the whole prize for the group. Yet the second attempt could
succeed when the first failed. They calculated both effects from the joint chance of success,
without assuming independence.

An early construction chose the best assignment among distinct jobs. Ada withdrew that restriction
after checking a move to an occupied job. Emmy independently showed that a collision can be the
better group choice. The peers then chose a target that maximises expected group energy across
\emph{all} job-or-idle profiles. They ranked its assigned agents ahead of visitors and gave the full
prize to the first successful agent. Emmy checked Ada's equilibrium argument independently; their
derivations are in \texttt{ada/priority\_prize\_note.md} and \texttt{emmy/research-note.md}.

Emmy then raised a survival objection: the same total prize can leave one agent unable to work
later. The peers also checked a budget bound for peer scoring when success chances are private.
Their final source check found no data in P that could test the priority rule.

\section{Results}

Consider two agents attempting job \(j\). Let \(R_j\) be its fixed prize, \(p_{1j}\) and \(p_{2j}\)
the agents' chances of success, and \(q_j\) the chance that both succeed. Under P's split rule, the
expected group payout is
\[
R_j(p_{1j}+p_{2j}-q_j).
\]
Agent 2's expected private prize is \(R_j(p_{2j}-q_j/2)\). Yet its added expected payout for the
group, with agent 1's attempt held fixed, is only \(R_j(p_{2j}-q_j)\). Its private incentive to join
the job therefore exceeds that contribution by \(R_jq_j/2\). The identity does not require
independent successes. It also does not say every collision is harmful: the second attempt still
adds the chance that agent 2 succeeds while agent 1 fails. Ada derived this comparison and Emmy
checked it in their notes.

\begin{proposition}
Fix one round in which active agents choose a job or remain idle. Suppose each agent maximises its
expected own energy, success and expected cost laws are known and stay fixed when other agents
change jobs, and there are no donations or later-round payoffs. Choose a job-or-idle profile that
maximises expected group energy over all feasible profiles. Rank agents assigned to each job in that
profile above visitors, and pay the job's entire prize to its highest-ranked successful agent. Then
the chosen profile is a Nash equilibrium of job choice. For every realised outcome, the rule pays
the same total job prize as P's split rule.
\end{proposition}

To see why, consider one agent leaving its assigned job. Its old priority prize is at least the
payout the group loses when it leaves: it can win even when a lower-ranked assigned agent also
succeeds. At a new job, the deviator ranks below those assigned there. It wins precisely when it
succeeds and every assigned agent at that job fails. Its new prize is therefore exactly the payout
it adds to the group at that job, including when the job was empty. The change in its own attempt
cost is also the change in group cost under the stated assumptions. Its gain from any job or idle
deviation is thus no greater than the change in expected group energy. The latter cannot be positive
at a global maximum. This is the argument Ada gave and Emmy checked; it allows correlated successes
and a target profile with collisions.

The rule only redirects a prize of size \(R_j\); adjustments to individual payouts sum to zero on
each job. This conditional result does not show that agents in P know the laws, follow
expected-energy best responses, or achieve the target.

\section{What did not hold, and what remains unsettled}

The distinct-job version of the construction did not hold in general. Maximising group energy only
among distinct assignments leaves open a profitable move into an occupied job. Allowing collisions
in the global optimisation settled that objection for the one-round game. It did not settle whether
the target can be found from P's published observations.

Nor does equal total payout mean equal survival. In Emmy's example, two agents each start with
\(15\) energy, spend \(10\) each, and both solve a job worth \(20\). Splitting leaves balances
\((15,15)\), while priority leaves \((25,5)\). A later call costing \(10\) can deactivate the
priority loser although both agents remain active after splitting. Donations, later work, and
changes in answer effort can therefore matter even when the current group total agrees. The
proposition does not address those effects.

The priority rule also assumes the planner knows success chances. Emmy's separate scoring
calculation shows why learning private chances may need a budget. Suppose two private types each
gain at least \(M\) from a false report before scoring, their distinct peer-outcome laws are \(P_0\)
and \(P_1\), and score payments lie between zero and \(B\). Here \(\operatorname{TV}(P_0,P_1)\) is
the total variation distance between those laws. The two truth-telling constraints imply
\[
B\geq \frac{M}{\operatorname{TV}(P_0,P_1)}.
\]
Ada checked the derivation. This bound applies to that payment model if scores consume P's energy.
It does not rule out every information mechanism, and Q's abstract rewards need not be P-energy.

Finally, P's reported collision counts cannot establish whether priority would improve group energy.
The conditions also differ in discussion cost, numbers of active agents, donations, and job mix. The
reports lack the counterfactual success chances and job-specific costs needed to value an
alternative assignment. The peers therefore made no measured improvement or novelty claim.

\section{Why the thread stopped}

The verifier accepted the one-round, prize-preserving equilibrium argument and its limits, but kept
the thread at DRAFT. The missing link is evidence about net energy and survival when real agents
choose, communicate, spend tokens, and face later rounds. The smallest useful restart is a
controlled comparison of P's split rule and the priority rule with the same jobs and agent
conditions, recording choices, joint successes, token costs, energy balances, deactivations, and
donations. Those records would test the incentive calculation and expose any survival cost of
reallocating the prize.

\bibliographystyle{plainurl}
\bibliography{references}

\end{document}
